% Drawing 75/01-2 -- folder 75, batch 01, archive page 5 (PDF page 6).
% Transcription: dans la marge gauche, deux dessins : au crayon, un contour à deux sommets et deux arêtes ; à l'encre, un pentagone
% Kind: contour with two vertices and two edges; pentagon. Method 2 (schematic redraw from the facsimile). Status: ok.
% Notes: A contour of two vertices joined by two edges (a lens), and a pentagon with marked vertices. On the render the lens is the heavier stroke and the pentagon the lighter one, the reverse of what the transcription says (pencil lens, ink pentagon); both are drawn as on the page.
% Coordinates: 1 unit = 1 mm = 1/100 of the crop width (bbox in page fractions: [0.1, 0.18, 0.36, 0.38]).
% Redrawn by Claude (claude-opus-5-5), issue #73, 2026-11-07; not checked by a human; not compiled (no TeX in the session).
% Rights: reproduces part of the work -- behind RENDER_TRANSCRIPTS, never published by itself.
\documentclass[tikz,border=3pt]{standalone}
\usepackage{amsmath,amssymb}
\usetikzlibrary{patterns,arrows.meta}
\tikzset{pencil/.style={black!45}, struck/.style={black!35}, mark/.style={circle,draw,inner sep=1.3pt}, pt/.style={circle,fill,inner sep=0.9pt}, opt/.style={circle,draw,fill=white,inner sep=1.1pt}}
\begin{document}
\begin{tikzpicture}[x=1mm,y=1mm,line cap=round,line join=round,font=\scriptsize]
\draw[very thick] plot[smooth, tension=0.55] coordinates {(10,-40) (13,-25) (33,-17)};
\draw[very thick] plot[smooth, tension=0.55] coordinates {(10,-40) (25,-36) (33,-17)};
\node[pt] at (10,-40) {};
\node[pt] at (33,-17) {};
\draw (50,-38.5) -- (74,-38) -- (82,-64) -- (46,-85) -- (33,-62) -- cycle;
\node[pt] at (74,-38) {};
\node[pt] at (82,-64) {};
\node[pt] at (46,-85) {};
\node[pt] at (33,-62) {};
\end{tikzpicture}
\end{document}
